\section{The Symplectic Presentation}\label{sec:symplectic}

\subsection{Divide and conquer}

Factor the group; present the factors; compose the presentations.  Two
short exact sequences organise the qupit Clifford group:
\[
  \begin{array}{l@{\qquad}l}
  1 \to \PPauli \to \PCliff \to \Sp \to 1 & \text{semidirect: it splits;}\\[2pt]
  1 \to \langle\omega\rangle \to \Cliff \to \PCliff \to 1 & \text{central extension.}
  \end{array}
\]
This section presents the symplectic factor, which is where the real
rewriting happens; \S\ref{sec:composing} presents the Pauli factor and
the scalars and composes.  The plan for $\Sp$ has seven steps, which the
subsections follow: the normal form; rewriting to it, warming up with
$S_n$; the passage from symmetric to symplectic; uniqueness; relation
reduction; the presentation theorem; and a comparison with the qubit
case.

\subsection{The semantics: symplectic transformations of Pauli vectors}\label{sec:spsem}

Rather than $2n \times 2n$ matrices, the formalisation models $\Sp$ as
the group of $\Zp$-linear, symplectic-form-preserving bijections of the
phase space $\aid{Pauli}\;n = \aid{Vec}\;(\Zp \times \Zp)\;n$ --- one
pair $(a,b)$, the exponents of $Z^{a}X^{b}$ (up to phase), per wire.
A symplectic transformation carries its own inverse and its own proofs:

\begin{agdacode}
record Symplectic (n : ℕ) : Set where
  field
    ap        : Pauli n → Pauli n
    ap⁻¹      : Pauli n → Pauli n
    invˡ      : ∀ p → ap⁻¹ (ap p) ≡ p
    invʳ      : ∀ p → ap (ap⁻¹ p) ≡ p
    linear-+  : ∀ p q → ap (p +ₚ q) ≡ ap p +ₚ ap q
    linear-*  : ∀ k p → ap (k *ₚ p) ≡ k *ₚ ap p
    preserves : ∀ p q → sform (ap p) (ap q) ≡ sform p q

_≈ˢ_ : ∀ {n} → Symplectic n → Symplectic n → Set
S ≈ˢ T = ap S ≗ ap T
\end{agdacode}

Equality is extensional equality of the action alone, so the proof
fields never obstruct.  Composition is function composition, and
therefore associativity and the unit laws hold \emph{on the nose}, by
\aid{refl} --- a small choice with large consequences downstream, since
every monoid case of every soundness proof closes definitionally.
\aid{Sp-group}~$n$ is the resulting \aid{Group}.  The generators act
as the qupit paper's conjugation tables (its Lemmas~2.22--2.23) read on
exponent vectors:

\begin{agdacode}
actg : ∀ {n} → Gen n → Pauli n → Pauli n
actg (gate₁ H-gate)  ((a , b) ∷ ps)             = (- b , a) ∷ ps
actg (gate₁ S-gate)  ((a , b) ∷ ps)             = (a , b + a) ∷ ps
actg (gate₂ CZ-gate) ((a , b) ∷ (a' , b') ∷ ps) = (a , b + a') ∷ (a' , b' + a) ∷ ps
actg (g ↥)           (x ∷ ps)                   = x ∷ actg g ps

⟦_⟧ : ∀ {n} → Circuit n → Symplectic n
⟦ [ g ]ʷ ⟧ = ⟦ g ⟧ᵍ
⟦ ε ⟧      = εˢ
⟦ w • v ⟧  = ⟦ w ⟧ ∘ˢ ⟦ v ⟧
\end{agdacode}

Words are read left to right, so $w \bullet v$ applies $w$ first and its
denotation composes on the right.  The seven proof obligations per
generator --- inverses, linearity, preservation of the form --- are
identities in $\Zp$, discharged with the ring solver.  Soundness of
every axiom of \S\ref{sec:reduction} in this semantics is one clause per
axiom in \texttt{SoundnessDirect} (1\,100 lines).

\subsection{The normal form: boxes, rows, and the staircase}\label{sec:nf}

The qupit paper's symplectic normal form is a $Z$-normal layer followed
by an $X$-normal layer on $n$ wires, then a normal form on $n-1$ wires.
We refine this into a doubly inductive datatype.  The atomic boxes were
shown in \S\ref{sec:background}; rows and the staircase are:

\begin{minipage}[c]{0.36\textwidth}
\begin{agdacode}[basicstyle=\footnotesize\ttfamily]
M L ML : ℕ → Set

M 0 = ⊤
M (₁₊ n) = Vec D n × E

L 0 = ⊤
L (₁₊ n) = Vec B n × A

ML 0 = ⊤
ML 1 = M 1 × L 1
ML m@(₂₊ n) =
  M m × L m ⊎ D × ML (₁₊ n)
\end{agdacode}
\end{minipage}\hfill
\begin{minipage}[c]{0.29\textwidth}\centering
  \begin{tabular}{@{}r@{\,}c@{\,}l@{}}
    \ufig[0.6]{m4-box} & \scalebox{0.7}{$=$} & \ufig[0.6]{m4-atoms}\\[1ex]
    \ufig[0.6]{l4-box} & \scalebox{0.7}{$=$} & \ufig[0.6]{l4-atoms}\\
  \end{tabular}
\end{minipage}\hfill
\begin{minipage}[c]{0.29\textwidth}\centering
  \begin{tabular}{@{}r@{\,}c@{\,}l@{}}
    \ufig[0.6]{ml4-box} & \scalebox{0.7}{$=$} & \ufig[0.6]{ml-tower-a-r}\\[1ex]
    \ufig[0.6]{ml4-box} & \scalebox{0.7}{$=$} & \ufig[0.6]{ml-tower-b-r}\\
  \end{tabular}
\end{minipage}

\medskip\noindent
An $M$-row is a column of $D$ boxes ending in an $E$ box (the paper's
$X$-normal layer); an $L$-row is a column of $B$ boxes ending in an $A$
box (its $Z$-normal layer, with the $A$ box at the bottom).  An
$\ML$ box at width $m \ge 2$ is \emph{either} a full pair of rows
\emph{or} a bottom $D$ box followed by an $\ML$ box one wire up: a
spine of $D$ boxes capped by a pair of rows.  This is the first
induction, on the width of a single coset representative.  The second
induction is the staircase:

\begin{minipage}[c]{0.54\textwidth}
\begin{agdacode}[basicstyle=\footnotesize\ttfamily]
NF : ℕ → Set
NF 0 = ⊤
NF (₁₊ n) = NF n × ML (₁₊ n)

[_] : ∀ {n} → NF n → Word (Gen n)
[_] {0}     tt         = ε
[_] {₁₊ n}  (nf , ml)  = [ nf ] ↑ • [ ml ]ᵐˡ
\end{agdacode}
\end{minipage}\hfill
\begin{minipage}[c]{0.42\textwidth}\centering
  \ufig[0.7]{nf4}
\end{minipage}

\medskip\noindent
Reading normal forms back as circuits is structural recursion, box by
box.  The box-to-word maps make the pictures of \S\ref{sec:background}
precise; for instance
\begin{agdacode}
[_]ᵈ : ∀ {n} → D → Word (Gen (₂₊ n))
[_]ᵈ {n} (₀ , b) = Ex • CZ^ (- b)
[_]ᵈ {n} (a@(₁₊ _) , b) = Ex • CZ^ (- a) • H • S^ -b/a
\end{agdacode}
where $-b/a = -b \cdot a^{-1}$, and an $A$ box $A_{a,b}$ with $a \ne 0$
is $\aid{XM}\;a \bullet H \bullet S^{-b/a}$, a multiplier, a Hadamard,
and an $S$-power.  Why the spine?  The paper's $Z$-normal circuit
$W_Z^{(n)}$ has its $A$ box on some wire $i$, with $B$ boxes above it
and nothing below; the position $i$ is an index that the paper's
normal-form circuits carry implicitly.  The $\ML$ type makes it explicit
as the length of the $D$-spine, and, crucially, makes the coset type at
width $m$ a \emph{sum} of the coset type at width $m-1$ and a flat case
--- exactly the shape the coset table and the uniqueness proof recurse
on.  The identity coset is the flat one with all exponents zero and
$A_{0,1}$ at the bottom, and the section sends it to $\varepsilon$ up to
$\approx$, because $\mathit{Ex}^2 \approx \varepsilon$ collapses each
$D_{0,0}$ and $B_{0,0}$.

\subsection{Warm-up: coset tables for the symmetric group}\label{sec:sn}

Swap circuits on $n$ wires present $S_n$, and the library's treatment of
them~\cite{qupit-code} is the template the symplectic development
follows.  Circuits have a built-in tower of subgroups
$\mathsf{Cir}\,0 < \mathsf{Cir}\,1 < \dots < \mathsf{Cir}\,n$, where
$\mathsf{Cir}\,k$ consists of the circuits that touch only the top $k$
wires; for swap circuits this is $S_0 < S_1 < \dots < S_n$.
Representatives of $S_3/S_2$ (and of $S_2/S_1$, $S_1/S_0$) are the
descending staircases:
\begin{center}
\begin{tabular}{ccc@{\qquad\qquad}cc@{\qquad\qquad}c}
  \ufig[0.8]{c2-0} & \ufig[0.8]{c2-1} & \ufig[0.8]{c2-2} &
  \ufig[0.8]{c1-0} & \ufig[0.8]{c1-1} & \ufig[0.8]{c0-0}\\
  $c^{(2)}_0$ & $c^{(2)}_1$ & $c^{(2)}_2$ & $c^{(1)}_0$ & $c^{(1)}_1$ & $c^{(0)}_0$
\end{tabular}
\end{center}
For $f \in S_3$ there is a unique $c \in C_2$ with $c \circ f^{-1}$
fixing the bottom wire --- the $c$ that sends $f^{-1}(0)$ to $0$ ---
and recursing on $c \circ f^{-1} \in S_2$ writes every $f$ uniquely as
$c^{(0)} c^{(1)} c^{(2)}$: the mixed-radix encoding $|S_3| = 3 \cdot 2
\cdot 1$ of \S\ref{sec:engine} (Figure~\ref{fig:pushing}, right).

Rewriting to normal form is the coset table.  Pushing a generator $b$
through a coset representative $c_3 \in C_3$ yields a new coset $c_3'$
and a residual ``dirty'' circuit $b'$ on the wires above, which is then
pushed through $C_2$, and so on:

\begin{minipage}[c]{0.62\textwidth}
\begin{agdacode}[basicstyle=\footnotesize\ttfamily]
ract : C (₁₊ n) → Gen (₂₊ n) → Circuit (₁₊ n) × C (₁₊ n)
ract {n}    ε         σ-gen  = ε , σ• ε
ract {n}    (σ• ε)    σ-gen  = ε , ε
ract {₁₊ n} (σ• σ• c) σ-gen  = (σ {n = n}) , σ• σ• c
ract {n}    ε         (g ↥)  = [ g ]ʷ , ε

ract-sound : ∀ {n} c b →
  let (b' , c') = ract {n} c b
  in  [ c ]ᶜ • [ b ]ʷ ≈ b' ↑ • [ c' ]ᶜ
\end{agdacode}
\end{minipage}\hfill
\begin{minipage}[c]{0.36\textwidth}\centering
  \ufig[0.55]{ract-l}\;\scalebox{0.7}{$\approx$}\;\ufig[0.55]{ract-r}
\end{minipage}

\medskip\noindent
The four cases are the four pictures of Figure~\ref{fig:pushing}:
disjoint wires commute; the Yang--Baxter rule passes the swap through
and shifts it up; $\sigma^2 = e$ shrinks the staircase; the staircase
grows.  Every case has the shape $c \cdot b \approx b' \cdot c'$, the
five coset-table hypotheses close by evaluation, and collecting the
cases gives a rewrite system that normalises every circuit.
Simplifying the collected relations leaves the two Coxeter rules,
$\sigma^2 = \varepsilon$ and Yang--Baxter, plus the structural ones, and
the library proves that they present the permutation group of
$\mathsf{Fin}\;n$.

\begin{figure}
\centering
\begin{minipage}[c]{0.66\textwidth}\footnotesize
\begin{tabular}{r@{\,}c@{\,}l@{\quad}l}
  \ufig[0.5]{push1-l} & \scalebox{0.7}{$\rightarrow$} & \ufig[0.5]{push1-r} &
    disjoint wires: commute past \\
  \ufig[0.5]{push2-l} & \scalebox{0.7}{$\rightarrow$} & \ufig[0.5]{push2-r} &
    Yang--Baxter: pass through, shift up \\
  \ufig[0.5]{push3-l} & \scalebox{0.7}{$\rightarrow$} & \ufig[0.5]{push3-r} &
    $\sigma^2 = e$: staircase shrinks \\
  \ufig[0.5]{push4-l} & \scalebox{0.7}{$\rightarrow$} & \ufig[0.5]{push4-r} &
    staircase grows \\
\end{tabular}
\end{minipage}\hfill
\begin{minipage}[c]{0.30\textwidth}\centering
  \ufig[0.7]{sn-stair}
\end{minipage}
\caption{Left: pushing a swap through a staircase, the four cases of the
coset table for $S_n$.  Right: the staircase of staircases, the normal
form of a permutation.}
\label{fig:pushing}
\end{figure}

\subsection{From symmetric to symplectic: change the action}\label{sec:pushing}

The same framework serves the symplectic group under a richer action.
There are three ways of looking at the wires.  $S_n$ acts on
$\{0,\dots,n-1\}$: a wire is a position, and $c \circ f^{-1}$ fixes
$0$.  $S_n$ acts on $\aid{Vec}\;\mathsf{Bit}\;n$: a wire is a bit, and
$c \circ f^{-1}$ fixes the $\mathbb{Z}_2$-linear subspace of the bottom
wire.  Adding $H$, $S$, $CZ$, the group acts on
$\aid{Vec}\;\mathrm{Pauli}_1\;n$, a $2n$-dimensional symplectic
$\Zp$-vector space in which a wire is a two-dimensional symplectic
subspace; enlarging the coset layer from staircases to $\ML$ boxes, for
each $f \in \Sp$ there is a $c$ such that $c \circ f^{-1}$ fixes the
symplectic subspace of the bottom wire.  The cosets at width $1{+}n$
are the $\ML$ boxes, their section is $[\_]^{\mathsf{ml}}$, and the coset
table is the paper's ``push a gate through a normal box'', now a
function:

\begin{agdacode}
C = ML

ract : ∀ {n} → C (₁₊ n) → Gen (₁₊ n) → Circuit n × C (₁₊ n)

ract-sound : ∀ {n} c g →
  let (b' , c') = ract {n} c g
  in [ c ]ᶜ • [ g ]ʷ ≈ b' ↑ • [ c' ]ᶜ
\end{agdacode}

\aid{ract} dispatches on the shape of the coset --- a flat box or a
$D$-spine --- and on the generator --- $S$ or $H$ on the bottom wire,
$CZ$ on the bottom two, or a lifted gate --- and delegates each case to
a module of the \texttt{Pushing} directory (26 files, 5\,600 lines):
$S$ and $H$ through a flat box, $CZ$ through a flat box, $CZ$ meeting a
$D$ box and a flat box, a lifted gate through the spine.  The residual
$b'$ lives one wire up, as the type says.  The soundness proof
\aid{ract-sound} is a chain of $\approx$-steps per clause, each step a
box relation or a structural rule.  For example, an $S$ on the bottom
wire meeting a spine $D_{a,b}$ followed by $\mathit{lm}$: $S$ commutes
past the lifted $\mathit{lm}$, then the box relation
$D \leftarrow S$ rewrites $D_{a,b} \bullet S$ to $\mathit{dir}{\uparrow}
\bullet D_{a,b-a}$, and re-association finishes.

So far this is the $S_n$ story with bigger cosets.  The difference is in
the five hypotheses of \S\ref{sec:engine}.  Four close as before; the
second, \aid{h-wd-ax} --- that the table respects the axioms, i.e.\
congruent inputs produce equal cosets and congruent residuals --- does
\emph{not} close by evaluation, because the axioms of $\Sp$ (unlike
$\sigma^2 = e$) are not decided by a finite computation on the coset
data.  The formalisation proves it in two halves.  The \emph{coset half}
is proved semantically: apply $\sem{\_}$ to \aid{ract-sound} on both
sides of an axiom, observe that residuals live one wire up and so cannot
move the bottom wire, and conclude from head-injectivity of the coset
section (\S\ref{sec:uniqueness}) that the two cosets are equal.  The
\emph{residual half} then follows by cancellation --- provided the shift
$\uparrow$ is injective on congruence classes at the lower width, which
is a consequence of faithfulness (completeness) there.  This makes the
tower an induction:

\begin{agdacode}[basicstyle=\footnotesize\ttfamily]
record TowerLevel (n : ℕ) : Set where
  field
    nfp'  : NormalForm (n QRel,_===_) (NFᵗ n)
    agree : ∀ u → PB._≈_ (n QRel,_===_) (SNF.NormalForm.inv-nf nfp' u) [ φ u ]

-- The induction.  Level (1+k) needs ↑-injectivity at k, which comes
-- from faithfulness at k, which comes from Level k -- strictly below,
-- so the recursion is well founded.
tower : ∀ n → TowerLevel n
tower 0 = record { nfp' = base0' ; agree = λ _ → PB.refl }
tower (₁₊ k) = record { nfp' = nfp'₁ ; agree = agree₁ }
  where
  prev = tower k
  ↑inj = SemInj.Injectivity!.↑-inj-↑ k
           (FF.faithful-from′ k (TowerLevel.nfp' prev) φ φ-inj (TowerLevel.agree prev))
  nfp'₁ = T.Extension.nfp' (ext k ↑inj) (TowerLevel.nfp' prev)
\end{agdacode}

The loop is: a normal form at width $k$ gives faithfulness at $k$
(\aid{faithful-from′}, by \aid{by-normalization}); faithfulness gives
injectivity of $\uparrow$ on congruence classes (\aid{SemInj}, using
that a lifted word acts as $\mathrm{id} \oplus \sem{w}$); that
discharges \aid{h-wd-ax} at width $k$; and the generic
\aid{Extension.nfp'} produces the normal form at width $k{+}1$.  The
base case, width $1$, is built by hand.  This is the formal counterpart
of the paper's Lemma~4.2 and Proposition~4.3: no termination argument is
needed beyond structural recursion, and ``the dirty gates eventually get
removed'' is the type of \aid{ract}.

\subsection{Uniqueness of the normal form}\label{sec:uniqueness}

Theorem~\ref{thm:unique} says that normal forms with equal denotations
are equal.  The action of a normal form on a Pauli vector peels the
staircase one wire at a time --- the coset factor acts, the resulting
bottom entry is fixed, and the inner normal form acts on the tail ---
and injectivity is a structural induction on the staircase:

\begin{agdacode}[basicstyle=\footnotesize\ttfamily]
act-nf : ∀ {n} → NF n → Pauli n → Pauli n
act-nf {0}    _        = id
act-nf {₁₊ n} (ih , lm) ps = head s ∷ act-nf ih (tail s)
  where s = act [ lm ]ᵐˡ ps

lemma-lm-head-inj : ∀ {n} (lm₁ lm₂ : ML (₁₊ n)) →
  (∀ (ps : Pauli (₁₊ n)) → head (act [ lm₁ ]ᵐˡ ps) ≡ head (act [ lm₂ ]ᵐˡ ps)) → lm₁ ≡ lm₂

lemma-lm-tail-surj : ∀ {n} (lm : ML (₁₊ n)) (qs : Pauli n) →
  ∃ λ (ps : Pauli (₁₊ n)) → tail (act [ lm ]ᵐˡ ps) ≡ qs

lemma-nf-inj : ∀ {n} (nf₁ nf₂ : NF n) → act-nf nf₁ ≗ act-nf nf₂ → nf₁ ≡ nf₂
\end{agdacode}

Two obligations.  \emph{Tail surjectivity} is free, because a
symplectic transformation carries its inverse.  \emph{Head injectivity}
--- the bottom-wire output determines the coset --- is the crux, and it
is itself an induction on the $\ML$ type (\texttt{LMHeadInj},
1\,300 lines): at width $1$ a case analysis on the four shapes of an
$A$ box using closed forms for the actions of $S^{k}$, $HS^{k}$ and the
multipliers; a separation lemma showing a flat box and a spine never
agree on heads; and, for two spines, the $D$ box is read off by probing
with the two basis Paulis $Z$ and $X$ on the bottom wire and the identity
above, after which the hypothesis descends one width.  With these two
lemmas the generic \texttt{Circuit.Uniqueness} module of the library
assembles \aid{⟦[]⟧-injective} and the \aid{UniqueNormalForm} witness;
the module's header records the one subtlety, that \aid{ap} is a
\emph{left} action --- \aid{ap} of a composite applies its right factor
first --- so the coset factor of a normal form is the one next to the
input, which is what lets the coset be recovered from every state rather
than from states of a restricted shape.  This is the paper's
Proposition~3.8 with its Lemmas~3.4--3.7, now an induction whose
structure is dictated by the datatype.

\subsection{Relation reduction: box relations and the two rule sets}\label{sec:reduction}

Which relations does \aid{ract-sound} use?  Rewriting \emph{every}
circuit to normal form and collecting every relation the rewrite ever
uses gives the paper's box relations --- one per gate per box shape,
parametrised by the box's labels.  The formalisation organises them as
nine families of case tables, each a lemma quantified over the box
parameters with a computed direction \aid{dir-of} and updated box:

\begin{center}\small
\begin{tabular}{@{}lll@{}}
\toprule
wires & family & pushed \\
\midrule
1 & $A \leftarrow H$, $A \leftarrow S$, $E \leftarrow S$ & $H$ or $S$ through an $A$ box; $S$ through an $E$ box \\
2 & $B \leftarrow H{\uparrow},S{\uparrow},S$; \ $D \leftarrow H,S{\uparrow},S,CZ$; \ $L \leftarrow CZ$ & a gate through a $B$, a $D$, or an $L$-row \\
3 & $BB \leftarrow CZ{\uparrow}$, $B{\uparrow} \leftarrow CZ$, $DD \leftarrow CZ$ & $CZ$ through two $B$ or two $D$ boxes \\
\bottomrule
\end{tabular}
\end{center}

For instance, the two-wire table for a $D$ box is

\begin{agdacode}
D←H-S↑-S-CZ : ∀ (d : D) (g : Gen 2) (neq : g ≢ H-gen ↥) →
  let (e , dir) = TD.dir-of d g neq ; d' = TD.d'-of d g neq
  in [ d ]ᵈ • [ g ]ʷ ≈ S^ e • dir ↑ • [ d' ]ᵈ

d'-of (a , b) H-gen _       = (b , - a)
d'-of (a , b) S-gen _       = (a , b + - a)
d'-of (a , b) (S-gen ↥) _   = (a , b)
d'-of (a , b) CZ-gen _      = (a , b + - ₁)
\end{agdacode}

--- the paper's Figure~18, with the residual split into an $S$-power on
the bottom wire and a word one wire up, and with the box update a
function of the labels.  The nine families are the paper's
Figures~15--21; unfolding the case tables gives 42 parametric relations,
against the paper's count of 66, the difference being only in how cases
are grouped (the formalisation, for instance, proves each of the
$\uparrow$/$\downarrow$ pairs of a $B$ box and a $D$ box once and
obtains the other by a duality lemma).  The width-3 relations are stated
at width~3 and widened to arbitrary width by a congruence
$\aid{cong↓ᵏ}$.

The box relations are derived (\texttt{BR}, 7\,000 lines, on top of a
calculus of $\mathit{Ex}$-conjugations and two-qupit lemmas of 11\,000
lines) from the \emph{full} symplectic rule set, seventeen group-specific
axioms whose two- and three-qupit members are Selinger's
c10--c15~\cite{selinger2015clifford} verbatim and whose one-qupit
members are the multiplier calculus:

\begin{agdacode}
order-S    : (₁₊ n) SRel,  S ^ p === ε
order-H    : (₁₊ n) SRel,  H ^ 4 === ε
order-SH   : (₁₊ n) SRel,  (S • H) ^ 3 === ε
comm-HHS   : (₁₊ n) SRel,  H • H • S === S • H • H
M-mul      : ∀ x y → (₁₊ n) SRel,  M x • M y === M (x *' y)
semi-MS    : ∀ x →   (₁₊ n) SRel,  M x • S === S^ (x ^2) • M x
semi-M↑CZ  : ∀ x →   (₂₊ n) SRel,  M x ↑ • CZ === CZ^ (x ^1) • M x ↑
semi-M↓CZ  : ∀ x →   (₂₊ n) SRel,  M x ↓ • CZ === CZ^ (x ^1) • M x ↓
\end{agdacode}

together with \aid{order-CZ}, \aid{comm-CZ-S↓}, \aid{comm-CZ-S↑} and
\aid{selinger-c10} to \aid{c15}.  Here $M_x$ is spelled
$S^{x}HS^{x^{-1}}HS^{x}H$, and four of the axioms are \emph{schemes}
over units $x, y$.  By inspection one finds a smaller set of primitive,
bidirectional relations that derive all of these: the \emph{simplified}
symplectic rule set, which replaces the schemes by statements about the
single multiplier $M_g$ of the primitive root, strengthens $H^4 =
\varepsilon$ to $H^2 = M_{-1}$, and drops \aid{order-SH} and
\aid{comm-HHS}:

\begin{agdacode}
order-S :           (₁₊ n) SRel,  S ^ p === ε
order-H :           (₁₊ n) SRel,  H ^ 2 === M₋₁
M-power :     ∀ k → (₁₊ n) SRel,  Mg^ k === M (g^ k)
semi-MS :           (₁₊ n) SRel,  Mg • S === S^ (g * g) • Mg
semi-M↑CZ :         (₂₊ n) SRel,  Mg ↑ • CZ === CZ^ g • Mg ↑
semi-M↓CZ :         (₂₊ n) SRel,  Mg ↓ • CZ === CZ^ g • Mg ↓
\end{agdacode}

plus the same \aid{order-CZ}, \aid{comm-CZ-S↓}, \aid{comm-CZ-S↑} and
\aid{selinger-c10} to \aid{c15}: fifteen group-specific rules, or
eighteen with the three structural ones, shown as circuits in
Figure~\ref{fig:simp}.  The choice is not canonical; one picks the set
one finds natural.  The two rule sets generate the same congruence, and
since they share an alphabet the identity on words is an isomorphism of
presentations in both directions (\texttt{Simplified.Iso}); the
substantive direction derives each full axiom from the simplified ones.
The multiplier law is typical:

\begin{agdacode}
M-mul : ∀ n (x y : ℤ* ₚ) → let open PB (Sim._QRel,_===_ (₁₊ n)) in
        M x • M y ≈ M (x *' y)
M-mul n = SimL.Lemmas1.lemma-M-mul n
\end{agdacode}

--- \fitfig[0.3\textwidth]{mm-stmt} --- write $x = g^{k}$, $y = g^{l}$,
use \aid{M-power} twice, add the exponents, reduce modulo $p-1$ using
$M_g^{\,p-1} \approx \varepsilon$, which is Fermat's little theorem
(\S\ref{sec:modp}), and use \aid{M-power} once more.  The whole
reduction from seventeen rules to fifteen is 900 lines; the reduction
from the box relations to the seventeen is 18\,000.  This is the
formal counterpart of the paper's Theorem~4.4, whose proof the paper
already delegated to Agda.

\begin{figure}
\centering\footnotesize
\begin{minipage}[t]{0.46\textwidth}\centering
  \begin{tabular}{@{}r@{\;\,}r@{\,}c@{\,}l@{}}
    order-S & \ufig[0.55]{simp-order-S-l} & $=$ & \ufig[0.55]{simp-order-S-r}\\[1pt]
    order-H & \ufig[0.55]{simp-order-H-l} & $=$ & \ufig[0.55]{simp-order-H-r}\\[1pt]
    M-power & \ufig[0.55]{simp-M-power-l} & $=$ & \ufig[0.55]{simp-M-power-r}\\[1pt]
    semi-MS & \ufig[0.55]{simp-semi-MS-l} & $=$ & \ufig[0.55]{simp-semi-MS-r}\\[1pt]
    order-CZ & \ufig[0.55]{simp-order-CZ-l} & $=$ & \ufig[0.55]{simp-order-CZ-r}\\
  \end{tabular}
\end{minipage}\hfill
\begin{minipage}[t]{0.52\textwidth}\centering
  \begin{tabular}{@{}r@{\;\,}r@{\,}c@{\,}l@{}}
    comm-CZ-S↓ & \ufig[0.55]{simp-comm-CZ-Sdn-l} & $=$ & \ufig[0.55]{simp-comm-CZ-Sdn-r}\\[1pt]
    comm-CZ-S↑ & \ufig[0.55]{simp-comm-CZ-Sup-l} & $=$ & \ufig[0.55]{simp-comm-CZ-Sup-r}\\[1pt]
    semi-M↑CZ & \ufig[0.55]{simp-semi-Mup-CZ-l} & $=$ & \ufig[0.55]{simp-semi-Mup-CZ-r}\\[1pt]
    semi-M↓CZ & \ufig[0.55]{simp-semi-Mdn-CZ-l} & $=$ & \ufig[0.55]{simp-semi-Mdn-CZ-r}\\
  \end{tabular}
\end{minipage}\\[3pt]
\begin{tabular}{@{}r@{\;\,}r@{\,}c@{\,}l@{}}
  selinger-c10 & \ufig[0.55]{simp-selinger-c10-l} & $=$ & \ufig[0.55]{simp-selinger-c10-r}\\[1pt]
  selinger-c11 & \ufig[0.55]{simp-selinger-c11-l} & $=$ & \ufig[0.55]{simp-selinger-c11-r}\\
\end{tabular}\\[3pt]
\begin{tabular}{@{}r@{\;\,}r@{\,}c@{\,}l@{}}
  selinger-c12 & \ufig[0.45]{simp-selinger-c12-l} & $=$ & \ufig[0.45]{simp-selinger-c12-r}\\[1pt]
  selinger-c13 & \ufig[0.45]{simp-selinger-c13-l} & $=$ & \ufig[0.45]{simp-selinger-c13-r}\\
\end{tabular}\\[1pt]
\begin{tabular}{@{}r@{\;\,}r@{\,}c@{\,}l@{}}
  selinger-c14 & \ufig[0.45]{simp-selinger-c14-l} & $=$ & \ufig[0.45]{simp-selinger-c14-r}\\[1pt]
  selinger-c15 & \ufig[0.45]{simp-selinger-c15-l} & $=$ & \ufig[0.45]{simp-selinger-c15-r}\\
\end{tabular}
\caption{The simplified symplectic rule set: fifteen group-specific rules
(here $M_g$ and $M_{-1}$ abbreviate the $SHS$-words of the text).}
\label{fig:simp}
\end{figure}

\subsection{The symplectic presentation theorem}\label{sec:sptheorem}

Soundness (one clause per axiom), the normal form of \S\ref{sec:pushing}
and its uniqueness (\S\ref{sec:uniqueness}) assemble, through
\aid{by-normalization}, into a sub-presentation: the interpretation is an
injective homomorphism from the words modulo the full rules into
\aid{Sp-group}.  Surjectivity --- every symplectic transformation is the
action of some circuit --- is proved constructively by ``clear the bottom
wire, recurse on the tail'': the single-level input is

\begin{agdacode}
Theorem-LM : ∀ {n} (p q : Pauli n) → sform p q ≡ ₁ →
  ∃ λ (lm : ML n) → act [ lm ]ᵐˡ p ≡ pZ₀ × act [ lm ]ᵐˡ q ≡ pX₀
\end{agdacode}

--- for any symplectic pair there is a coset box sending it to the
bottom-wire $Z$ and $X$ --- which is the paper's Lemma~3.7 with the box
computed rather than asserted.  Composing the sub-presentation of the
full rules with the identity isomorphism of \S\ref{sec:reduction} gives
the theorem for the simplified rules, with the \emph{same}
interpretation:

\begin{agdacode}
Theorem-Sp : ∀ {n} → (n SRel,_===_) IsPresentationOf (Sp-group n)
Theorem-Sp = SimPres.presentation
\end{agdacode}

\subsection{Comparison with the qubit case}

Selinger's Figure~8~\cite{selinger2015clifford} presents the qubit
Clifford group with the scalars and Paulis included.  Read modulo Paulis
and scalars, its three-qubit relations are \emph{identical} to
\aid{selinger-c12} to \aid{c15} --- the reason those names were kept ---
and its two-qubit relations are \emph{similar}, with $S^{-1}$ for $S$
and $CZ^{p-1}$ for $CZ$.  The one-qupit relations are \emph{more
complicated}: the odd-prime multiplier calculus of $M_g$, with its
dependence on a primitive root and on Fermat's theorem, has no qubit
counterpart, where the only multiplier is $M_{-1} = H^2$ up to phase.
That the difficulty of the general odd prime concentrates on one wire is
what makes the divide-and-conquer route affordable: the two- and
three-wire rewriting is Selinger's, and the one-wire arithmetic is
arithmetic.
